\title{Very-high-frequency oscillations in the main peak of a magnetar giant flare}
\begin{document}
\maketitle

\begin{abstract}
Nature | Vol 600 | 23/30 December 2021 | 621 Article Very-high-frequency oscillations in the main peak of a magnetar giant flare A. J. Castro-Tirado1,2, N. Østgaard3 ✉, E. Göǧüş4 ✉, C. Sánchez-Gil5, J. Pascual-Granado1, V. Reglero6,7, A. Mezentsev3 ✉, M. Gabler6 ✉, M. Marisaldi3,8 ✉, T. Neubert9, C. Budtz-Jørgensen9, A. Lindanger3, D. Sarria3, I. Kuvvetli9, P. Cerdá-Durán6, J. Navarro-González7, J. A. Font6,10, B.-B. Zhang11,12,13, N. Lund9, C. A. Oxborrow9, S. Brandt9, M. D. Caballero-García1, I. M. Carrasco-García14, A. Castellón2,15, M. A. Castro Tirado1,16, F. Christiansen9, C. J. Eyles7, E. Fernández-García1, G. Genov3, S. Guziy17,18, Y.-D. Hu1,19, A. Nicuesa Guelbenzu20, S. B. Pandey21, Z.-K. Peng11,12, C. Pérez del Pulgar2, A. J. Reina Terol2, E. Rodríguez1, R. Sánchez-Ramírez22, T. Sun1,23,24, K. Ullaland3 \& S. Yang3 Magnetars are strongly magnetized, isolated neutron stars1–3 with magnetic fields up to around 1015 gauss, luminosities of approximately 1031–1036 ergs per second and rotation periods of about 0.3–12.0 s. Very energetic giant flares from galactic magnetars (peak luminosities of 1044–1047 ergs per second, lasting approximately 0.1 s) have been detected in hard X-rays and soft γ-rays4, and only one has been detected from outside our galaxy5. During such giant flares, quasi-periodic oscillations (QPOs) with low (less than 150 hertz) and high (greater than 500 hertz) frequencies have been observed6–9, but their statistical significance has been questioned10. 
\end{abstract}

\section{Recovered Text}
Nature  |  Vol 600  |  23/30 December 2021  |  621
Article
Very-high-frequency oscillations in the main
peak of a magnetar giant flare
A. J. Castro-Tirado1,2, N. Østgaard3 ✉, E. Göǧüş4 ✉, C. Sánchez-Gil5, J. Pascual-Granado1,
V. Reglero6,7, A. Mezentsev3 ✉, M. Gabler6 ✉, M. Marisaldi3,8 ✉, T. Neubert9,
C. Budtz-Jørgensen9, A. Lindanger3, D. Sarria3, I. Kuvvetli9, P. Cerdá-Durán6,
J. Navarro-González7, J. A. Font6,10, B.-B. Zhang11,12,13, N. Lund9, C. A. Oxborrow9, S. Brandt9,
M. D. Caballero-García1, I. M. Carrasco-García14, A. Castellón2,15, M. A. Castro Tirado1,16,
F. Christiansen9, C. J. Eyles7, E. Fernández-García1, G. Genov3, S. Guziy17,18, Y.-D. Hu1,19,
A. Nicuesa Guelbenzu20, S. B. Pandey21, Z.-K. Peng11,12, C. Pérez del Pulgar2, A. J. Reina Terol2,
E. Rodríguez1, R. Sánchez-Ramírez22, T. Sun1,23,24, K. Ullaland3 \& S. Yang3
Magnetars are strongly magnetized, isolated neutron stars1–3 with magnetic fields
up to around 1015 gauss, luminosities of approximately 1031–1036 ergs per second
and rotation periods of about 0.3–12.0 s. Very energetic giant flares from galactic
magnetars (peak luminosities of 1044–1047 ergs per second, lasting approximately 0.1 s)
have been detected in hard X-rays and soft γ-rays4, and only one has been detected
from outside our galaxy5. During such giant flares, quasi-periodic oscillations (QPOs)
with low (less than 150 hertz) and high (greater than 500 hertz) frequencies have been
observed6–9, but their statistical significance has been questioned10. High-frequency
QPOs have been seen only during the tail phase of the flare9. Here we report the
observation of two broad QPOs at approximately 2,132 hertz and 4,250 hertz in
the main peak of a giant γ-ray flare11 in the direction of the NGC 253 galaxy12–17,
disappearing after 3.5 milliseconds. The flare was detected on 15 April 2020 by the
Atmosphere–Space Interactions Monitor instrument18,19 aboard the International
Space Station, which was the only instrument that recorded the main burst phase
(0.8–3.2 milliseconds) in the full energy range (50 × 103 to 40 × 106 electronvolts)
without suffering from saturation effects such as deadtime and pile-up. Along with
sudden spectral variations, these extremely high-frequency oscillations in the burst
peak are a crucial component that will aid our understanding of magnetar giant flares.
We report here the detection11 of a new giant flare (initially dubbed
GRB 200415) with the Atmosphere–Space Interactions Monitor
(ASIM) aboard the International Space Station (ISS) on 15 April 2020
at 08:48:05.56 (±0.03) ut. With the ASIM Modular X- and Gamma-ray
Sensor (MXGS) instrument18,19, we recorded data for 2 s centred around
the burst. The two independent detectors of MXGS, covering energies
between 50–400 keV (low-energy detector, LED) and 300 keV to 40 MeV
(high-energy detector, HED), did not suffer from saturation effects
(deadtime, pile-up) and recorded for the first time the fine structure
of the main burst phase (0.8–3.2 ms) of a magnetar in this entire energy
range (Fig. 1, Extended Data Fig. 2). Owing to the large effective area of
ASIM, 1-μs time resolution and large energy range, we have performed
both a detailed time analysis and a spectral analysis of the main phases
of the giant flare. We are able to resolve the complex temporal structure
prior to the absolute peak emission, consisting of six distinct intensity
peaks during the first 3.2 ms, this flare being the first one for which we
have seen multiple peaks prior to the maximum (see Fig. 1). During
the approximately 160-ms duration of the giant flare, around 1046 erg
isotropic equivalent energy was released, roughly the energy the Sun
radiates in about 100,000 yr.
https://doi.org/10.1038/s41586-021-04101-1
Received: 17 August 2020
Accepted: 6 October 2021
Published online: 22 December 2021
 Check for updates
1Instituto de Astrofísica de Andalucía del Consejo Superior de Investigaciones Científicas (IAA-CSIC), Granada, Spain. 2Unidad Asociada al CSIC Departamento de Ingeniería de Sistemas y
Automática, Escuela de Ingeniería Industrial, Universidad de Málaga, Málaga, Spain. 3Birkeland Centre for Space Science, Department of Physics and Technology, University of Bergen, Bergen,
Norway. 4Faculty of Engineering and Natural Sciences, Sabancı University, Istanbul, Turkey. 5Departamento de Estadística e Investigación Operativa, Universidad de Cádiz, Puerto Real, Spain.
6Departamento de Astronomía y Astrofísica, Universitat de València, Valencia, Spain. 7Imaging Processing Laboratory, Universitat de València, Valencia, Spain. 8National Institute for
Astrophysics, Osservatorio di Astrofisica e Scienzia dello Spazio, Bologna, Italy. 9DTU Space, Technical University of Denmark, Kongens Lyngby, Denmark. 10Observatori Astronòmic, Universitat
de València, Valencia, Spain. 11School of Astronomy and Space Science, Nanjing University, Nanjing, China. 12Key Laboratory of Modern Astronomy and Astrophysics, Nanjing University,
Ministry of Education, Nanjing, China. 13Department of Physics and Astronomy, University of Nevada, Las Vegas, Las Vegas, NV, USA. 14Departamento de Química Analítica, Universidad de
Málaga, Málaga, Spain. 15Departamento de Álgebra, Geometría y Topología, Facultad de Ciencias, Universidad de Málaga, Málaga, Spain. 16Escuela Técnica Superior de Arquitectura,
Universidad de Málaga, Málaga, Spain. 17Astronomical Observatory, Mykolaiv National University, Mykolaiv, Ukraine. 18Research Institute, Mykolaiv Astronomical Observatory, Mykolaiv, Ukraine.
19Facultad de Ciencias, Universidad de Granada, Granada, Spain. 20Thüringer Landessternwarte Tautenburg, Tautenburg, Germany. 21Aryabhatta Research Institute of Observational Sciences
(ARIES), Nainital, India. 22Istituto di Astrofisica e Planetologia Spaziali, INAF, Rome, Italy. 23Purple Mountain Observatory, Chinese Academy of Sciences, Nanjing, China. 24School of Astronomy
and Space Science, University of Science and Technology of China, Hefei, China. ✉e-mail: Nikolai.Ostgaard@uib.no; ersin.gogus@sabanciuniv.edu; Andrey.Mezentsev@uib.no; michael.
gabler@uv.es; Martino.Marisaldi@uib.no
622  |  Nature  |  Vol 600  |  23/30 December 2021
Article
On the basis of the temporal structure and energy distribution we
distinguish four distinct phases: (i) precursor (0–0.8 ms, with flux
dominated by low-energy counts); (ii) burst (0.8–3.2 ms, with well
defined peak structure); (iii) decay (3.2–8.0 ms); and (iv) tail (8–160 ms).
For all four distinct phases we created detailed energy spectral fits
(see Methods for details), which are summarized in Fig. 2, Table 1 and
Methods (Extended Data Fig. 3). The preferred values for all of the four
phases are marked with bold text in Table 1.
The spectrum of the precursor phase is dominated by a blackbody
with a temperature of 66 keV. During the peak interval (0.8–3.2 ms)
the flux increases by one order of magnitude and we see a major
non-thermal component that can be modelled by a power law with
exponential cut-off with peak energy of 1,160 keV (maximum of the
spectral energy distribution). The flux of this non-thermal component
then gradually decreases during the decay and tail phase. In the tail
phase the blackbody component reappears, now with a slightly lower
temperature (40 keV) than in the precursor phase.
To search for quasi-periodic oscillations in the signal, we cre-
ated individual light curves each with a time resolution of 50 μs for
the LED and HED data, with four different search windows: 0–5 ms,
0–10 ms, 0–50 ms and 0–100 ms. The resulting power spectral density
(PSD) reaches a Nyquist frequency of 104 Hz, which is high enough to
include the shortest dynamical timescales for a typical neutron star.
The Leahy-normalized20 PSD constructed using the LED data reveals
two statistically significant QPOs at frequencies 2,156 ± 150 Hz and
4,256 ± 323 Hz (Fig. 3a, Extended Data Fig. 5, Extended Data Table 2,
marked with bold). These oscillations are seen exclusively during the
first few milliseconds, which includes the main burst phase (0.8–3.2 ms;
see Fig. 1c, Fig. 3b, and Methods for the details of the PSD modelling),
and they vanish 3.5 ms after the burst onset (Extended Data Fig. 8c).
The detection of the QPO coincides with the increase of the non-thermal
emission. Because 4,256 Hz is twice that of the lower frequency oscilla-
tion within errors, it is probably the first harmonic of the broad feature
at 2,156 Hz. The signal-to-noise ratio (SNR) at 2,156 Hz corresponds
to a P value <10−3 (see Extended Data Table 2 for additional details).
The PSD constructed using the HED data also shows these QPOs at
2,185 ± 52 Hz and 4,200 ± 37 Hz (Extended Data Table 2, Extended Data
Fig. 6), and both are in agreement with the frequencies from the LED
data within errors. Additionally, the HED power spectrum indicates the
presence of another statistically significant oscillation at approximately
1,487 ± 33 Hz, marked with bold in Extended Data Table 2.
We have also searched the unbinned LED and HED event arrival times
for these signals using Z m
2  statistics, where m is the number of har­
monics over which the Fourier powers are added21. We selected m = 1,
for which Z2 values correspond to twice the powers of the Rayleigh test.
We applied this technique in a dynamic manner. In particular, we con-
sidered a search interval of 5 ms, moving in steps of 1 ms with the start
time ranging from −6 ms through to 50 ms with respect to the burst
trigger time. In the frequency domain, we searched in 1,000 steps within
a 600-Hz frequency interval centred on the above-reported oscillation
GRB 200415A
103
–100
0
100
200
300
400
500
600
LED
HED
HED
LED
a
b
c
d
Counts per 5 ms
102
101
Time (ms)
103
–100
0
–2
0
2
4
6
8
10
30
20
10
0
30
20
10
0
100
200
300
400
500
600
Counts per 5 ms
Counts per 0.05 ms
102
101
Time (ms)
Time (ms)
–2
0
2
4
6
8
10
Time (ms)
Fig. 1 | Temporal variability of GRB 200415A. a, The ASIM LED (low-energy detector)
hard X-ray light curve of the magnetar giant flare. b, The ASIM HED (high-energy
detector) γ-ray light curves of the magnetar giant flare. Both light curves have a
complex structure, displaying six intensity peaks during the first 3.2 ms.
The time interval ranges from −100 to +600 ms, with the data binned to 5 ms.
c, d, A magnification of the first 10 ms with the LED (c) and HED (d) data binned to 50 μs.
Table 1 | Best-fitting parameters for the spectral properties of the different time epochs
Phase
Time interval (ms)
Modela
Photon index
Peak energy (keV)
Temperature (keV)
Reduced χ2 (d.o.f.)
Fluxd (×10−5 erg cm−2 s−1)
Precursor
0–0.8
CPL
−
+
0.4 1.1
1.5
−
+
281 59
84
-b
0.50 (11)c
−
+
10.2 3.6
5.6
BB
-b
-b
66 12
15
−
+
0.52 (12)c
9.7 5.4
12
−
+
Peak
0.8–3.2
CPL
−
−
+
0.89 0.11
0.12
1,160 100
120
−
+
-b
1.49 (24)
−
+
85.7 7.1
7.7
Decay
3.2–8.0
CPL
−
−
+
0.41 0.15
0.16
−
+
1,400 100
110
-b
1.01 (24)
−
+
44.2 4.6
5.1
Tail
8.0–160
CPL
0.85 0.14
0.15
−
−
+
−
+
990 110
140
-b
1.96 (21)
−
+
2.26 0.15
0.16
BB + CPL
−
+
6.3 1.0
3.0
−
+
737 32
34
39.6 3.0
3.3
−
+
0.71 (19)
2.19 0.19
0.20
−
+
aCPL, cut-off power law; BB, blackbody.
cC statistics are used instead of χ2 owing to the low number of counts available.
dIn the 60 keV–10 MeV range, 68\% confidence.
Nature  |  Vol 600  |  23/30 December 2021  |  623
frequencies. For the LED, we obtained the highest Z2 power at
f1 = 2,132 Hz (Extended Data Table 3) in the time interval from 0 to 5 ms
(see Fig. 3b), both consistent with the PSD results. We determined the
chance probability of this signal as 2.4 × 10−9 (Extended Data Table 3).
A Z 1
2 search at even higher frequencies yields a peak power at
f2 = 4,250 Hz with a chance probability of 1.7 × 10−4 in the same time
interval. This strengthens the evidence that indicates that the highest
frequency signal is most probably the harmonic. Similarly, we have
performed Z 1
2 searches in exactly the same manner using the HED data.
We obtain clear signals at 2,095 Hz with a chance probability of
5.0 × 10−8, and at 4,127 Hz with a probability of arising from a random
data sample as 1.1 × 10−2. We also conducted another search centred on
1,400 Hz, evident from the PSD modelling of the HED data. We find a
significant Z 1
2 peak at 1,353 Hz with a chance probability of 1.2 × 10−12
in HED but much less significant (probability of approximately 0.05)
in LED (see Extended Data Table 3). This Article analyses only the QPOs
significant in both HED and LED, and in the following we will refer to
these frequencies, f1 = 2,132 Hz and f2 = 4,250 Hz, as identified in
LED data.
We have also investigated unweighted Swift/Burst Alert Telescope
(BAT) Gamma-ray Urgent Archiver for Novel Opportunities (GUANO)22–24
observations for the presence of high-frequency QPOs. Owing to the
internal data readout time limitation of 0.1 ms, we could construct the
light curve on this timescale, which provides a frequency coverage up
to 5,000 Hz. The Leahy-normalized Swift/BAT PSD clearly shows broad
signals with the peak power in the range 2,163–2,272 Hz (see Fig. 3a),
and other signals with peak power in the range 1,464–1,614 Hz, consist-
ent with all features we see in the ASIM-LED PSD. Swift/BAT GUANO
observations, therefore, provide an independent confirmation of the
f1 = 2,132 Hz QPO seen with ASIM-LED.
Variations on timescales t ≤ 0.5 ms can be caused by magnetic field
instabilities in the magnetosphere close to the stellar surface (stellar
radii, r ≤ 100 km). Independently of whether we consider a crustquake
or magnetospheric instabilities as the initial trigger of GRB 200415A,
Alfvén waves will be created in the magnetosphere. The Alfvén crossing
time, tA ≈ πr/c ≈ 1 ms (c, speed of light in a vacuum), of waves constrained
to this region determines the timescale of reconnection and is the typical
timescale during which instabilities develop25,26. While bouncing back and
forth between the footpoints of their respective magnetic field lines, these
waves interact with each other nonlinearly across different field lines,
dissipating energy. These ‘encounters’ produce emissions with a typical
variability on a timescale of 1 ms, which is consistent with f1 at 2,132 Hz.
This process could be very inefficient and last for tens of Alfvén crossing
times27. f2 = 4,250 Hz would then naturally be the overtone (within errors)
of the broad feature at f1, or could be produced by the interaction of multi-
ple Alfvén waves propagating along neighbouring field lines. After several
crossings most of the energy of the waves has been absorbed by the crust.
The reconnection is expected to terminate after a few milliseconds25,26,
which would naturally explain the sudden disappearance of the QPOs.
Furthermore, the observed timescale of t ≈ 1/f1 = 0.469 ms deter-
mines a lengthscale in which the energy of the giant flare is released,
l = c × t ≈ 140 km. Assuming for simplicity a spherical magnetic field
line, its radius would then be approximately 20 km. Hence the total
volume of the flare is expected to be comparable in magnitude to, or
somewhat larger than, the neutron star itself.
The properties of the observed spectra (Fig. 3, Table 1) are consistent
with what is expected to happen during a reconnection event in the
magnetosphere. The photospheric temperature of an expanding
fireball, which forms as a consequence of the energy release during
the reconnection, is expected to have typical temperatures in the
range 20–80 keV (ref. 25), consistent with the 40–60 keV observed.
The non-thermal component could be created by resonant cyclotron
scattering, which dominates the scattering optical depth in highly
magnetized magnetospheres and which can increase the energy
of thermally emitted photons by orders of magnitudes. The strong
magnetic field also enables us to understand the observed peak energy
a
102
10–2
10–4
10–6
100
Counts s–1 keV–1
b
 QF(Q) (erg cm–2 s–1)
10–4
10–6
102
10–2
100
c
10–4
10–6
102
10–2
100
d
10–4
10–6
102
10–2
100
102
103
104
102
103
104
Energy (keV)
Energy (keV)
102
103
104
102
103
104
Energy (keV)
Energy (keV)
102
103
104
102
103
104
Energy (keV)
Energy (keV)
102
103
104
102
103
104
Energy (keV)
Energy (keV)
Fig. 2 | Time-resolved spectroscopy. Spectra for four time intervals during
the magnetar giant flare, defined in the text. a, Precursor: 0.0–0.8 ms. b, Peak:
0.8–3.2 ms. c, Decay: 3.2–8 ms. d, Tail: 8–160 ms. A total of 2,148 counts in the
LED and 2,141 counts in the HED have been used. Left, count spectra (blue, LED;
red, HED) and best-fit model (green). Right, energy flux density for the best-fit
models listed in Table 1.
624  |  Nature  |  Vol 600  |  23/30 December 2021
Article
around 1 MeV, because in very strong magnetic fields processes, such as
one-photon pair creation and photon splitting, the energy of photons
is limited to roughly 2mec2, where me is the mass of an electron.
An alternative explanation of the timing features is based on the prox-
imity of the QPO candidate at f1 = 2,132 Hz to one of the high-frequency
QPOs observed in the tail of SGR 1806-20 with f = 1,840 Hz (ref. 9).
The high-frequency QPOs in magnetars are commonly interpreted as
radial overtones of the fundamental (magneto-)elastic oscillations
with one or more nodes in the crust, which will be preferably excited
during the flare28. Depending on how exactly the instability in the mag-
netosphere is triggered, there may be strong perturbations in the crust
of the neutron star which should naturally excite oscillations. Follow-
ing this interpretation, the second strong feature at f2 ≈ 4,250 Hz may
then be related to an even higher overtone. Within this interpretation
f1 can be considered as an upper estimate on the purely shear mode
with two (or three) nodes in the crust (see Methods section ‘QPO
theoretical implications’). This interpretation is also consistent with
theoretical expectations and the constraints29 that were obtained
from the QPOs of SGR 1806-20. However, the sudden disappearance
of the QPOs after approximately 3.5 ms and the evolution of energy
spectra slightly favour our first model, but do not exclude the presence
of stellar oscillations.
Online content
Any methods, additional references, Nature Research reporting sum-
maries, source data, extended data, supplementary information,
acknowledgements, peer review information; details of author con-
tributions and competing interests; and statements of data and code
availability are available at https://doi.org/10.1038/s41586-021-04101-1.
1.
Kaspi, V. M. \& Beloborodov, A. M. Magnetars. Annu. Rev. Astron. Astrophys. 55, 261–301 (2017).
2.
Mereghetti, S., Pons, J. A. \& Melatos, A. Magnetars: properties, origin and evolution.
Space Sci. Rev. 191, 315–338 (2015).
3.
Duncan, R. C. \& Thompson, C. Formation of very strongly magnetized neutron stars:
implications for gamma-ray bursts. Astrophys. J. 392, L9–L13 (1992).
4.
Mazets, E. P., Golenetskii, S. V., Il’inskii, V. N., Aptekar’, R. L. \& Guryan, Yu. A. Observations
of a flaring X-ray pulsar in Dorado. Nature 282, 587–589 (1979).
5.
Ofek, E. O. et al. The short-hard GRB 051103: observations and implications for its nature.
Astrophys. J. 652, 507–511 (2006).
6.
Barat, C. et al. Fine time structure in the 1979 March 5 gamma ray burst. Astron. Astrophys.
126, 400–402 (1983).
7.
Strohmayer, T. E. \& Watts, A. L. Discovery of fast X-ray oscillations during the 1998 giant
flare from SGR 1900+14. Astrophys. J. Lett. 632, L111–L114 (2005).
8.
Israel, G. L. et al. The discovery of rapid X-ray oscillations in the tail of the SGR 1806-20
hyperflare. Astrophys. J. 628, L53–L56 (2005).
9.
Strohmayer, T. E. \& Watts, A. L. The 2004 hyper flare from SGR 1806-20: further evidence
for global torsional vibrations. Astrophys. J. 653, 593–601 (2006).
10.
Pumpe, D., Gabler, M., Steininger, T. \& Enβlin, T. A. Search for quasi-periodic signals in
magnetar giant flares – Bayesian inspection of SGR 1806-20 and SGR 1900+14. Astron.
Astrophys. 610, A61 (2018).
11.
Marisaldi, M., Mezentsev, A., Østgaard, N., Reglero, V. \& Neubert, T. GRB 200415A: ASIM
observation. GCN Circular 27622 (2020).
12.
Svinkin, D. et al. A bright γ-ray flare interpreted as a giant magnetar flare in NGC 253.
Nature 589, 211–213 (2021).
13.
Roberts, O. R. et al. Rapid variability of a giant flare from a magnetar in NGC 253. Nature
589, 207–210 (2021).
14.
The Fermi-LAT Collaboration. High-energy emission from a magnetar giant flare in the
Sculptor galaxy. Nat. Astron. 5, 385–391 (2021).
15.
Yang, J. et al. GRB 200415A: a short gamma-ray burst from a magnetar giant flare?
Astrophys. J. 899, 106 (2020).
16.
Zhang, H. M., Liu, R.-Y., Zhong, S.-Q. \& Wang, X.-Y. Magnetar giant flare origin for GRB
200415A inferred from a new scaling relation. Astrophys. J. Lett. 903, 32 (2020).
17.
Burns, E. et al. Identification of a local sample of gamma-ray bursts consistent with a
magnetar giant flare origin. Astrophys. J. Lett. 907, 28 (2021).
18.
Neubert, T. et al. The ASIM mission on the International Space Station. Space Sci. Rev.
215, 26 (2019).
19.
Østgaard, N. et al. The Modular X- and Gamma-Ray Sensor (MXGS) of the ASIM payload
on the International Space Station. Space Sci. Rev. 215, 23 (2019).
20.	 Leahy, D. A. et al. On searches for pulsed emission with application to four globular cluster
X-ray sources: NGC 1851, 6441, 6624 and 6712. Astrophys. J. 266, L160–L170 (1983).
21.
Buccheri, R. et al. Search for pulsed γ-ray emission from radio pulsars in the COS-B data.
Astron. Astrophys. 128, 245–251 (1983).
22.	 Gehrels, N. et al. The Swift gamma-ray burst mission. AAstrophys. J. 611, 1005 (2004).
23.	 Barthelmy, S. D. et al. The Burst Alert Telescope on the Swift Midex mission. Space Sci.
Rev. 120, 143–164 (2005).
24.	 Tohuvavohu, A. et al. Gamma-Ray Urgent Archiver for Novel Opportunities (GUANO):
Swift/BAT event data dumps on demand to enable sensitive subthreshold GRB searches.
AAstrophys. J. 900, 35 (2020).
25.	 Mahlmann, J. F., Akgün, T., Pons, J. A., Aloy, M. A. \& Cerdá-Durán, P. Instability of twisted
magnetar magnetospheres. Mon. Not. R. Astron. Soc. 490, 4858–4876 (2019).
26.	 Parfrey, K., Beloborodov, A. M. \& Hui, L. Twisting, reconnecting magnetospheres and
magnetar spindown. Astrophys. J. Lett. 754, 12 (2012).
27.
Li, X., Zrake, J. \& Beloborodov, A. M. Dissipation of Alfvén waves in relativistic
magnetospheres of magnetars. Astrophys. J. 881, 13 (2019).
28.	 Duncan, R. C. Global seismic oscillations in soft gamma repeaters. Astrophys. J. 498,
L45–L49 (1998).
29.	 Gabler, M., Cerdá-Durán, P., Stergioulas, N., Font, J. A. \& Müller, E. Constraining properties
of high-density matter in neutron stars with magneto-elastic oscillations. Mon. Not. R.
Astron. Soc. 476, 4199–4212 (2018).
Publisher’s note Springer Nature remains neutral with regard to jurisdictional claims in
published maps and institutional affiliations.
© The Author(s), under exclusive licence to Springer Nature Limited 2021
Fig. 3 | Periodogram and fits for quasi-periodicities search for the time
interval 0–5 ms. a, Periodograms for ASIM-LED (50–400 keV; red) and Swift/
BAT GUANO (15–150 keV; black) observations. The interval 0–200 ms is used for
both periodograms. ASIM time resolution is 50 μs (10-kHz upper frequency) and
Swift/BAT time resolution is 100 μs (5-kHz upper frequency). The blue horizontal
line is the white noise level. The power spectral densities show broad signals after
700 Hz, providing independent confirmation that the f1 = 2,132 Hz QPO
determined from the ASIM-LED data is genuine. b, Dynamic power contours
(in red) resulting from the Z2 search in the 1,800–2,400 Hz frequency range,
along with the LED light curve with 50-μs time resolution. The inset is the
expanded view of the Z2 contours corresponding to the 99\% (innermost),
95\% (middle) and 90\% (outermost) levels of peak Z2 power centred at 2,132 Hz
(Extended Data Table 3). For comparison the frequency found by the PSD analysis
at 2,156 ± 45 Hz (the frequency and its uncertainty) is shown by dashed and
dotted horizontal lines (Extended Data Table 2, time interval 0–10 ms).
100
101
2,400
2,200
2,150
2,100
2,050
1.0
1.5
2.0
2.5
2,200
Time (ms)
Time (ms)
2,000
1,800
–2
0
2
4
6
8
10
102
103
25
Intensity (counts per 50 μs)
20
15
10
5
0
104
Frequency (Hz)
Frequency (Hz)
Frequency (Hz)
104
103
102
101
100
Leahy normalized PSD
Swift-BAT
LED
White noise level
b
a
Methods
Instrumentation and data acquisition
Instrument description and observation geometry. The Modular X-
and Gamma-ray Sensor (MXGS) is an imaging and spectroscopic X-ray
and γ-ray instrument19 mounted on the starboard side of the Columbus
module on the International Space Station. Together with the Modu-
lar Multi-Spectral Imaging Assembly (MMIA)30 MXGS constitutes the
instruments of the Atmosphere–Space Interactions Monitor (ASIM).
The main objectives of MXGS are to image and measure the spectrum of
X-rays and γ-rays from lightning discharges, known as terrestrial γ-ray
flashes (TGFs). MMIA is designed to image and perform high-speed pho-
tometry of transient luminous events (TLEs) and lightning discharges.
With these two instruments specifically designed to explore the relation
between electrical discharges, TLEs and TGFs, ASIM is the first mission
of its kind. Owing to telemetry limitations, ASIM is a triggered system;
for each trigger, 2 s of data centred at the event is downloaded.
At the magnetar giant flare time, the ISS was 434.74 km above the
Pacific, close to the New Zealand coast, at longitude 176.287° E and
latitude 39.442° S. As the flare was localized6 by the Inter Planetary
Network (IPN) on a region (a 4.73 arcmin × 3.78 arcmin parallelogram
centred at coordinates RAJ2000: 11.874°, dec.J2000: −25.194°) consistent
with the NGC 253 galaxy (Extended Data Fig. 1), this was 71.25° off-axis,
outside the MXGS imaging field of view (FOV). This is very close to the
Earth limb, which covers 69.41° in the MXGS FOV, but it is high enough
to rule out Earth’s atmospheric absorption effects. Using the GEANT4
Monte Carlo simulation toolkit, considering the minimum altitude
the photons reached (73.9 km) and the total column density traversed
(1.57 × 10−3 g cm−2) towards the direction of the NGC 253 galaxy, we see
no effect on photons above 50 keV in terms of time delays or energy
spectrum distortion, as the mean free path for 50-keV photons for this
column density is >106 km.
Data processing and validation. From both HED and LED we have 2 s
of data centred on the time of GRB200415A, with temporal resolution
of 1 µs (LED) and 28 ns (HED). For the analysis in this paper energies be-
tween 50 keV and 400 keV are used for LED (geometric area 1,024 cm2)
and >300 keV for HED (geometric area 900 cm2). Although MXGS was
designed for a different purpose, it provided unprecedented observa-
tions of the GRB 200415A magnetar flare, both regarding temporal
resolution and spectral coverage. ASIM MXGS is the only instrument
that recorded the main burst phase in the energy range from 50 keV
up to 40 MeV without suffering from saturation effects (deadtime,
pile-up). Here we describe the performance and limitations of the two
instruments and emphasize that during the observations of the burst
the observed flux and the measured energies were not affected by
these limitations.
Instrument performance. It should be emphasized that both HED
and LED are two independent instruments with independent read-out
electronics. As for all instruments, HED and LED will miss counts when
the flux exceeds a certain level. During GRB 200415A the flux never
reached such a high level and here we explain why.
HED has 12 independent BGO/PMT bars. When an event occurs less
than 10 µs after the previous one in a single BGO/PMT bar, the second
event will make a pulse that sits on the tail of the previous pulse, and
is flagged as a 'fast' event. During the most intense part (0–3.6 ms) of
GRB 200415A HED only recorded three 'fast' events out of a total of
393 raw counts.
LED has 16 independent readout chains and each chain has a dead-
time of 1.4 µs. When there is a hit in the same chain within 1.4 µs the two
hits will be recorded as a multi-hit event with one time tag. The pixel
address will be lost and the two signals will be added, but the number
of hits within the 1.4 µs will be recorded. Hits in different chains will be
recorded separately, but can have the same time tag if they occur within
the same 1 µs. During the most intense part of GRB 200415A (0–3.6 ms)
LED recorded 787 raw counts, of which 155 (19.7\%) were multi-hits or had
the same time tag. We use the interval 0–3.6 ms to ensure we include
the spike at 3.4–3.6 ms (see Fig. 1b). If this fraction of 19.7\% is due to
Compton scattering in the detector, detector assembly or structures
around the detector, they should only be counted as single hits for the
purpose of timing analysis.
The ASIM mass model, developed using GEANT4, includes a detailed
representation of the detectors as well as all surrounding structures and
is well suited for testing how many Compton events should be expected.
Incoming photons during the 0–3.6 ms interval with the spectral distri-
bution given in Table 1 and shown in Fig. 3 were simulated and entered
into the mass model from the correct incoming direction. It was found
that 11–13\% of these photons would deposit more than one count in the
LED, and consequently should only be counted as single hits.
We also performed a probability analysis for the same interval
(0–3.6 ms) to test how many real counts we could have missed owing
to the 1.4-µs deadtime in each single chain. The raw data were binned
in 30-µs bins, and the occurrence rate in each of the 120 bins (0–3.6 ms)
was used to estimate how many real counts we will miss owing to the
1.4-µs deadtime in each chain for each time bin. This analysis showed
that we might have missed a total of about 14 counts in this interval, and
they would have been recorded as multi-hits. If we consider that 14 of
the multi-hits should have been counted as two real hits and Compton
multi-hits reduced accordingly, there would be 17.6\% Compton hits
(instead of 19.7\%) in this time interval, which is a slightly larger frac-
tion than the expected 11–13\% determined by using the mass model.
Accepted data. It is not possible to distinguish the possible real counts
(among the multi-hits or with same time-tag in LED) from Compton
scattered hits, and so we have excluded all the counts that are pos-
sible Compton photons from the analysis (except for spectroscopy,
where raw counts are used). These Compton counts will also result in a
non-Poissonian background distribution. To remove these counts, we
proceeded as follows: Counts that are separated by less than 200 ns in
HED are only counted as one. Counts that have the same time tag in LED
(that is, within the same µs) as well as multi-hits are also counted as only
one count. This procedure will ensure that all Compton events will only
be counted as single events, and also any cosmic showers or particle
cascades in the 2 s of data will only be counted as single counts. After
applying this procedure, the white noise (background) has a Poissonian
distribution, and applying Leahy normalization20 for the PSDs for both
LED and HED indicates that the white noise has the expected constant
power of 2. For the most intense part of the burst (0–3.6 ms) we end up
with 363 and 663 accepted counts in HED and LED, respectively. For all
plots in this paper, except for the spectroscopy analysis and the spectra
shown in Fig. 3 and Extended Data Fig. 3, only accepted data are used.
ASIM was the only instrument that detected the peak intensity of the
GRB in the energy range from 50 keV to tens of MeV: In Extended Data
Fig. 2 we show GRB 200415A as measured by Fermi, ASIM and Swift/BAT.
In the two upper panels it can be seen that the peak of GRB 200415A
between 1.6 ms and 3 ms was clearly detected by both HED and LED,
but not by the Fermi/GBM BGO or NaI. The lower panel shows that ASIM
LED and Swift/BAT both detected the peak of GRB 200415A in energies
below 350 keV, but ASIM is the only instrument that recorded the main
burst in the high energy range (300 keV to 40 MeV). In Methods sec-
tion ‘QPO analysis’, it will be shown that it is indeed in the rising phase
of GRB 200415A that specific frequencies were observed by both LED
and HED, with a statistically significant power level above the white
noise level.
Time lag analysis. Following common practice for GRB studies we also
tested whether there is an energy dependence in the arrival times of
photons, especially during the initial hard spike emission of the giant
flare. A cross-correlation analysis for the interval −2 ms to 10 ms using
Article
50 µs bins and steps showed that there is no time lag between the two
light curves. Time lag (HED before LED) can only be seen if we narrow
the bin size and steps down to 4 μs or less.
Spectroscopy
Spectral analysis. To study the spectral evolution of the burst, we
divided the time series into four intervals, as mentioned in the main
text. For each time interval, we extracted the spectral files for LED and
HED. For LED we used all counts, excluding multi-hits events, for which
the energy information is accumulated over all simultaneous counts.
For HED we used all counts except FAST events (see instrument perfor-
mance) and all counts that are closer to the previous one in the same
detector than a specific ‘safety time’. This time depends on the detector
and the energy of the previous count, and ranges from 2 μs to about
20 μs for previous counts with energy above 20 MeV. This procedure
has been introduced to ensure the highest reliability of the HED energy
measurements, which are affected to some extent by the signal of the
previous count. The procedure has been developed for the spectral
analysis of terrestrial γ-ray flashes (TGFs), which exhibit fluxes much
higher than those recorded for this event and therefore suffer much
more from this effect. This procedure removes 53 out of a total of 2,141
counts above 400 keV (2.5\%) and therefore does not substantially im-
pact the spectral analysis. We do not use accepted data for the spectral
analysis because the Compton events, counted as one by the accepted
data procedure, are properly accounted for in the detector response
matrices (DRM) and therefore are relevant for an accurate spectral
analysis. We generated the DRM of the HED and LED detectors for the
incoming directions of the burst. DRMs are built using the full mass
model of the payload and the surrounding structures (the Columbus
Module on the ISS). We then use XSPEC v12.10.1f, a widely used gener-
al-purpose X-ray spectral-fitting programme that is part of the XANADU
package, for spectral fitting using the standard forward-folding ap-
proach. We fit all the time intervals in the energy range 60–350 keV for
LED and 0.4–20 MeV for HED. Background is accumulated over the
time interval −700 ms to −20 ms. Data were rebinned in order to have
a minimum 1σ significance above the background in each bin, or a
maximum of ten channels rebinned together. This latest condition is
applied only in the highest energy part of the HED spectra. Rebinning
is implemented on data files using the grppha utility from the FTOOLS
package (https://heasarc.gsfc.nasa.gov/docs/software/ftools/ftools\_
menu.html). We test different models, including a power law with ex-
ponential cut-off (CPL), a simple power law (PO), and a blackbody
spectrum (BB). The CPL is a custom model of the form
f E
E
( ) ≈
exp[ −
]
α
E α
E
( + 2)
p
 where α is the photon index and the peak en-
ergy Ep is the maximum of the spectral energy distribution. All the in-
tervals can be fitted with a CPL. The best-fit parameters are shown in
Table 1, and the fitted energy spectra and the corresponding residuals
are shown in Extended Data Fig. 3. The precursor can also be well fitted
with a blackbody spectrum with temperature 66
keV
−12
+15
, resulting in a
reduced χ2 comparable to that of the CPL model but only one free pa-
rameter; therefore, this model should be preferred. The preferred
values for all the four phases are marked with bold in Table 1.
The spectral evolution can be summarized as follows: The spectrum
of the precursor phase reveals a blackbody temperature of 66 keV.
During the peak phase (0.8–3.2 ms) the flux increases by an order of
magnitude and the spectrum is dominated by an accelerated compo-
nent with a peak energy of 1,160 keV. It is during this acceleration phase
that the two significant QPOs appear. The accelerated component
dominates also in the decay phase, but the flux has decreased. In the
tail phase the blackbody component reappears with a slightly lower
temperature (40 keV) than in the precursor phase and the accelerated
component is weaker.
Comparison with Fermi GBM. For comparison and validation pur-
poses we compared the ASIM MXGS spectra with the Fermi GBM results
provided in ref. 13. We produced the LED and HED spectra for the same
time intervals used in the Fermi analysis, and fitted the datasets with a
cut-off power law model as done for Fermi. The time intervals are simi-
lar, but not identical to, those chosen in this paper. Extended Data Fig. 4
shows the ASIM unfolded data and the 1σ confidence interval model
spectra for Fermi presented in Fig. 1d of ref. 13. Note that the unfolded
data depend on the observed data, the response of the instrument and
the considered theoretical best-fit model. Extended Data Table 1 reports
the spectral parameters and the calculated fluxes for both ASIM MXGS
and Fermi GBM. We note an overall good agreement between ASIM and
Fermi spectral parameters. In particular, the peak energy is compatible,
within errors, for both instruments for all time intervals, supporting the
goodness of the calibration and spectral performance of ASIM MXGS.
However, the calculated fluxes, although well compatible for intervals
1 and 4, are substantially different for intervals 2 and 3, ASIM reporting
a much larger flux than Fermi. We note that these are the time intervals
with the highest flux. In ref. 13 it is reported that interval 2 is affected by
saturation and includes a correction factor for the flux obtained by
comparing the GBM and Swift/BAT flux in the same time interval. We
suggest that this correction may be underestimated, and saturation
may be affecting GBM data in interval 3 as well. Conversely, we con-
firm that saturation in the GBM data does not significantly affect the
overall spectral shape, as reported in ref. 13. We also note, in interval 4,
a significant discrepancy between LED data and the GBM spectrum at
energies below 300 keV. A simple cut-off power law fit does not yield
a very good result, suggesting the need for an extra component at low
energies that we model with a blackbody spectrum.
QPO analysis
Methodology. We approach the problem of detecting periodic or
almost-periodic signals in noisy time series within a Bayesian frame-
work31,32 in modelling the observed periodogram by assuming any
low-frequency broadband variability to be due to a noise process.
Owing to the lack of knowledge of the burst emission mechanism, we
make a simple but conservative assumption that all broadband power in
the periodogram is supplied by a red (aperiodic) noise process in the form
of a power law. As opposed to QPO features—assumed to be fairly narrow
features on top of this process—our assumption will cause weak signals
at low frequencies to be buried in the higher variance of the broadband
noise. In return, it would yield a very low false positive detection rate33.
As in ref. 31, the periodogram refers to the squared Fourier transform
of the data, and it is assumed to be a sample of the underlying physical
process (the burst envelope and one or several noise processes).
We make use of the Stingray Python Package33 to perform this analysis.
We computed light curves and periodograms for LED and HED, for
several time segments after the onset of the burst: 0–5 ms, 0–10 ms,
0–50 ms and 0–100 ms. A larger time interval gives much better fre-
quency resolution. In each case, we produced a light curve by binning
the data to a time resolution of 1/2νNyquist = 50 µs, corresponding to a
Nyquist frequency of νNyquist ≈ 104 Hz. We chose this time resolution
since, on the basis of typical neutron star sizes34, we do not expect any
QPO above 10 kHz.
We search both the unbinned periodogram as well as the same peri-
odogram binned (to 3 and 5 times the original frequency resolution
of 106 Hz). Additionally, we smooth the spectra with a Wiener filter
with different smoothing factors (3 and 5) and compare results of the
search of binned periodograms with searches across the smoothed
periodograms.
An advantage of such a Bayesian approach is to provide a statistically
rigorous framework to test whether additional model components
(such as Lorentzian QPOs) are required by the data in the search for
periodicities and quasi-periodicities. The Bayesian procedure has three
main parts:
1. Broadband noise model selection. We address this task as a
model-selection problem, within a fully Bayesian analysis, to find the
preferred broadband noise model to represent the periodogram’s
low-frequency part between two nested models. The first and simpler
model, our null hypothesis H0, is just a power law to account for the red
noise, plus a constant to account for the Poisson noise in the detection
process. The second and more complex model, the alternative hypoth-
esis H1, is a broken power law (see equation 8 in ref. 32).
We used uniform prior probability distributions for all model param-
eters since no a priori knowledge of the magnetar giant flare emission is
possible. We obtained the joint posterior distribution of given param-
eters using Bayes’ theorem, taking the maximum a posteriori (MAP)
estimates of the model parameters as our choice for the parameter
estimation. We then performed a suit of Markov chain Monte Carlo
simulations (MCMCs) within the Bayesian framework to ensure that
the dataset is adequately described by the null hypothesis.
We compute the likelihood ratio test (LRT) statistic between the two
models, or TLRT, an often-used statistic for nested models, based on the
ratio of the likelihood maxima values. We use 500 MCMC ensemble walk-
ers with 100 samples each, after a burn-in phase with 200 samples for
each walker. Then we can generate predictive (simulated) periodogram
data from TLRT, a 1,000 sub-sample of the parameter set, and compute a
distribution for the statistic TLRT to compare with observed TLRT
obs to obtain
the P value. Given the distribution of TLRT, we can compute the correspond-
ing tail area probability, or the probability of obtaining a value of the test
statistic as extreme as the one observed, under the assumption of the
null hypothesis. The simulated data are fitted with both models H0 and
H1, following the same procedure performed on the observed data.
In all cases the preferred model is the simple power law, and we adopt
model H0 for the rest of our analysis of this burst.
2. Searching for (quasi-)periodicities. We search the periodogram
for the highest data/model outlier and compare this outlier to those
distributed by pure broadband noise to find narrow features that may
be candidates for a possible QPO (see Extended Data Figs. 5b, 6b).
We use the chosen broadband noise model fit to the periodogram
in the previous step to draw a new posterior predictive sample of the
parameter sets, and simulate 1,000 periodograms. In the same fash-
ion as outlined in the noise model selection, from these simulated
periodograms (all following the chosen noise model), we compute
the distribution of two other statistics, generated from the sample of
simulations, to compare the observed value of the statistic and compute
the corresponding P values.
One statistic is the summed-square residuals TSSE, which is an indica-
tor for how well the model fits the data. The second statistic is the
maximum ratio of observed to model power, T
R
= max ˆ
j
j
R
, where
R
I
S
ˆ = 2 /
j
j
j, and Ij and Sj are the observed and model powers. This is a
sensible statistic with which to investigate the narrow features35.
The factor of two normalizes the residuals in such a way that Rˆj will be
distributed as χ2 (see ref. 35).
We search for the highest data/model outlier, TR in the unbinned,
binned and smoothed periodograms (see results in Extended Data
Table 2). We can see how the features around 2,000–2,200 Hz and
4,000–4,400 Hz appear clearly and constantly. The few cases where
the maximum ratio is detected at higher frequencies correspond with a
higher frequency resolution, and with a smoothness factor of 3 or 5. So
that re-binning the PSD to 3 or 5 times the original resolution is enough
to detect these main features again. Because 4,256 Hz is consistently
twice that of the former periodicity, within errors, it is probably the
first harmonic of the broad feature at 2,156 Hz (both frequencies are
marked with bold in Extended Data Table 2, LED 0–100 ms).
3. Search for QPOs in the binned data using the identical approach
as for the model selection in the first step.  We approach the search
for QPOs in the data by following the same procedure applied to the
broadband noise model selection. Thus from a model-selection point
of view, and assuming a quasi-periodicity as another type of random
process, we fit the periodogram simultaneously with a broadband noise
process as well as one or more Lorentzians representing the QPO, as
the alternative hypothesis (H1), against the broadband noise model
only, which is the null hypothesis (H0).
We use the residuals of the data divided by the preferred broadband
noise model as a starting point from which to choose the alternative
hypothesis. This is to decide whether to fit one or more Lorentzians
together to model the broadband noise process.
Then we repeat the same procedure as above. We generate a large
number of simulated periodograms from the MCMC sample of the
posterior distribution of the parameter set, under the assumption of
the null hypothesis (only noise). We compute then the LRT statistic
from each simulated periodogram in turn, and finally compare the
observed value of the statistic to the distribution generated from the
sample of simulations, to get the corresponding P value associated with
this model-selection problem; see Extended Data Table 2.
Many P values for the LRT test are much smaller than 10−6. To limit the
computational costs, they were estimated by the upper tail probability
of the simulated posterior distributions, given the observed LRT, and
corrected by the bin frequency size for each time interval and binning.
Thus, the threshold for significance is adjusted for the number of dif-
ferent segments searched and the different binning.
Even so, the resulting P values (see Extended Data Table 2) are very
small. The observed reduction in the likelihood ratio between null and
the alternative hypothesis is larger than can be expected by chance if
the null hypothesis (data were only broadband noise) were true, and
we can reject it. However, this test cannot be seen as direct evidence
that the alternative hypothesis is true31.
Moreover, to check the practical significance, we compute the effect
size in terms of the discrepancy between the 95th percentiles of the sim-
ulated LRT distribution and the observed one (since only the upper tails
of the distributions are of interest). This is ε = |LRT
−LRT
|/LRT
.
obs
95
sim
obs
This effect size was always larger than 0.85—that is, a relative distance in
the percentile at least in a decrease of 15\%, which agrees with the P value
computation. The original data have a much higher discrepancy value or
LRT statistic than any replicated datasets under the null hypothesis.
Another possible effect size could be comparing the fractional
root-mean-square (RMS) amplitude in the QPO candidates, com-
pared with the RMS in the simulated periodograms, under the
broadband noise null hypothesis, over the same frequencies. This is
ε = |RMSobs − RMSsim|/RMSobs. The effect size is larger than 0.6 for the
components around 4,250 and 2,132 Hz for LED data—that is, a relative
distance in the simulated noised components RMS in a decrease of
40\% (or an observed RMS 2.5 times the simulated). In HED, simulated
RMS is over the 60\% decrease (or an observed RMS above 66\%). These
effect sizes combined with the estimated P values support the exist-
ence and significance of these QPO candidates around the 2,132 and
4,250 Hz frequencies.
On the Z2 search and significance of oscillations. We have also in-
vestigated the time evolution of the detected QPO features by employ-
ing the Z2 statistics21. For this purpose, we used unbinned event arrival
times and selected the number of harmonics as one—which is equiva-
lent to the Rayleigh statistics—to search for a periodic signal37. We first
performed the search in the broad frequency range from 300 Hz to
5,000 Hz with the frequency resolution of 1 Hz using both LED and HED
data collected within the first 5-ms time interval. We present the result-
ing Z 1
2 power spectra in Extended Data Fig. 7. Owing to the computa-
tional costs, we limited the finer resolution and dynamic search within
600-Hz frequency windows centred around the frequencies of
detected QPOs from the PSD analysis. We performed the Z2 search for
each frequency with the data of LED and HED individually but within
the same time intervals: we searched a time segment of 5 ms in the time
interval ranging from T0 − 6 ms to T0 + 50 ms, by moving the search
segment by 1-ms time steps.
Our Z 1
2 search with the LED data around the frequency 2,156 Hz
resulted in the highest power at a frequency of 2,132 Hz. We calculated
Article
the significance of this signal as follows. Note that the powers of Z2 with
one degree of harmonic follows a χ2 distribution with two degrees of
freedom. We first calculated the random occurrence (chance) prob-
ability of the highest Z 1
2 power for a single trial case. We then computed
the joint probability of having a Z2 power spectrum of 1,000 frequency
bins, which is the number of frequency steps used in our search and
serves as the number of trial frequencies. We find that the 2,132 Hz
signal is significant beyond the 99.9\% confidence level both in LED and
HED (Extended Data Fig. 7).
We calculated the chance occurrence probability for each
quasi-periodic signal seen with the LED and HED data following the
same method, and present them in Extended Data Table 3.
Duration of the QPOs. To investigate the duration of the high-frequency
QPOs we inspected periodograms of the GRB region within ±300 ms
from the GRB onset. A time window of length 200 ms was used to pro-
duce Leahy-normalized PSD periodograms, which were combined
into a spectrogram with time step of 1 ms, presented in Extended Data
Fig. 8a. Time axis corresponds to the left edge of the 200-ms long slid-
ing window. The sharp edge of the GRB spectral contents can be seen at
around −200 ms, when the right edge of the sliding window just enters
the GRB region. Another sharp edge at around 0 ms marks the moment
when the left edge of the sliding window starts leaving the GRB region.
Red noise is clearly seen up to 1 kHz. Before and after these sharp edges
only white noise is seen, except for frequencies <100 Hz in the tail of
the GRB (lasting for ~100 ms).
On the magnified segments of the spectrogram focused on fre-
quency range of 1 kHz to 5 kHz, and produced with a higher time reso-
lution (10-µs time step), the appearance and disappearance of the
high-frequency QPOs can be seen (Extended Data Fig. 8b, c). A clear
transition region with strong perturbations is seen between 1 ms and
3.5 ms, where the main fine time structure of the burst phase is con-
centrated. The motion of the sliding window across this area shows
the spectral evolution of the signal contents. After ~3.5 ms the 2,132-Hz
QPO smears out down to the white noise level, and we will refer to 3.5 ms
as the end of the QPO.
QPO theoretical implications
As discussed above, our analysis indicates two statistically significant
QPOs, one at f1 = 2,132 Hz and the other at f2 = 4,250 Hz: both seen only
during the peak phase (0–3.2 ms). Therefore, in this section, we focus
on the theoretical interpretation of these two frequencies. There are
two distinct physical processes that may result in signals with frequen-
cies in this range: oscillations of the neutron star excited during the
giant flare and magnetic reconnection events in the magnetosphere.
Neutron star oscillations. There are a number of different oscilla-
tion modes in neutron stars that may give rise to frequencies above
1 kHz (ref. 34): pressure modes, fundamental modes, space-time
modes, superfluid modes and magneto-elastic modes. Of these, the
latter have the strongest potential to create an observable temporal
variation of the electromagnetic emissions during a magnetar giant
flare28,38–40. Additionally, they are expected to be preferentially excited
during the flare28. The strongest support in favour of this explanation
is the detection of this kind of QPO in the decaying tail of previously
observed giant flares7–9,41–43 of SGR 1806-20 and SGR 1900+14, although
some doubt was cast on the significance of these detections recently10.
The QPOs can be divided into two groups with low (f ≲ 150 Hz) and
high (f ≳ 500 Hz) frequencies. In particular, the 1,840-Hz oscillation
detected9 in SGR 1806-20 is very close to f1. The low-frequency QPOs are
usually related to magneto-elastic oscillations without nodes (n = 0) in
the radial direction inside the crust, whereas the high-frequency QPOs
could be oscillations that have one or more nodes (n ≥ 1) inside the crust.
To estimate the properties of the neutron star we have used the
empirical formula for the different oscillation frequencies given in
ref. 29. If the QPOs indeed represent stellar oscillations, the high fre-
quencies indicate that they are connected to the overtones of the
magneto-elastic oscillations that have one or more modes in the crust,
ltn, where l indicates the angular dependency and n is the number of
nodes inside the crust in the radial direction. The dependence of these
oscillations on l is much weaker than the dependence on n, and we
therefore concentrate on the latter and refer to the l = 2 QPO from now
on. Even for the most extreme (most massive) neutron star models,
none of the equations of state considered in the literature provides
crustal shear mode frequencies for 2t1 of the order of 2 kHz or more36,44.
 The frequencies of the n = 1 modes are predicted44 to be approxi-
mately f ≈ 400–800 Hz, with very extreme cases from 200 up to
1,500 Hz. Therefore, it is reasonable to associate f1 with a higher n > 1
overtone. The presence of a signature at ~1,400 Hz indicates that this
may be associated with lower n > 0 modes. However, owing to its low
significance we will not consider it further and we will discuss this else-
where in more detail.
The second strong feature at f2 may then be related to a higher over-
tone. If the interpretation in terms of neutron star oscillations is cor-
rect, this would limit the available choices of the crust model, because
the order of the different overtones of the crustal shear modes is not
simply 1:2:3:4:5, and thus to match f1 and f2 requires some fine-tuning.
However, in reality the magnetic field complicates this analysis and
can increase the frequencies substantially29. Note that the maximal
frequency increase allowed when including magnetic fields limits the
frequency of 2t1 to about twice the frequency of the pure shear mode.
Otherwise the magnetic field would dominate and no high-frequency
QPO would be expected29. Thus the magnetic field helps to increase
the frequency, but cannot bring typical shear mode frequencies of 2t1
up to more than 2 kHz. Therefore, f1 can be considered as an upper
estimate on the purely shear mode frequency of the 2t2 QPO (in the
case that the equation of state predicts that f
f
> 2
t
1
2
2 , the conclusion
would be that f1 is an upper estimate on f t3
2 ). Without independent
estimates for the magnetic field or the properties of the star, the
parameter space spanned by the equation of state, mass of the neutron
star model, and magnetic field strength and configuration is too big
to further constrain the properties of the neutron star.
Reconnection in the magnetosphere. Although the frequency
f1 = 2,132 Hz is intriguingly similar to the highest frequencies of the
QPOs observed in SGR 1806-20, there is a fundamental difference
in both observations. For SGR 1806-20 in particular the 150-Hz and
625-Hz QPOs were observed for several tens of seconds during the
X-ray tail9—however, the data are consistent with short-lived signals
<0.5 s appearing and dissapearing43—whereas here the QPO candidate
is only observed for ~3 ms during the initial short peak. The timescale
t ≤ 1 ms associated with f1 = 2,132 Hz corresponds very well to the Alfvén
crossing time tA of the nearby magnetosphere, which is essentially
equal to the light crossing time there. For the Alfvén waves confined
to a region close to the stellar surface (r ≤ 100 km) this means tA ≈ 1 ms.
The Alfvén crossing time determines the timescale of reconnection
and is the typical timescale during which magnetospheric instabili-
ties develop25,26,45,46. Therefore, it is tempting to assume that the fast
variability of the emission on a timescale of 1 ms can be explained by
Alfvén waves dissipating energy. Recently it has been suggested that
dissipation in the magnetosphere owing to nonlinear interactions of
Alfvén waves could be an inefficient process and could last for tens of
Alfvén crossing times27,46. Most of the energy in these waves is absorbed
by the crust and re-emitted on longer timescales (weeks to months).
In this picture, the initial trigger of GRB 200415A (which could be either
a crustquake or a magnetospheric instability) generates Alfvén waves
in the magnetosphere, which propagate along the magnetic field lines.
These Alfvén waves bounce back and forth between the footpoints of
the field lines, naturally defining the Alfvén crossing time as the typi-
cal timescale. When Alfvén waves travelling along different field lines
interact nonlinearly, reconnection of the field lines occurs and energy
is released. These encounters produce emission with a typical vari-
ability on a timescale of 1 ms. The frequency f2 = 4,250 Hz would then
naturally be the overtone, or could be produced by the interaction of
multiple Alfvén waves propagating along the neighbouring field lines.
After approximately ten crossings most of the energy in the waves has
been absorbed by the crust. The longer-lasting emission probably has
a different origin, related to the evaporation of plasma trapped in the
magnetosphere by the ultra-strong magnetic field47.
Data availability
The ASIM data used for this study are available on the Zenodo repository,
https://doi.org/10.5281/zenodo.5242975. The Swift 100-μs-resolution
data are publicly available at the Swift BAT GUANO repository: https://
www.swift.psu.edu/guano. The Fermi NaI and BGO data are publicly
available at https://heasarc.gsfc.nasa.gov/W3Browse/fermi/fermig-
brst.html.
Code availability
For the search for QPOs within the Bayesian framework, we use the
Stingray Python Package developed by Huppenkothen et al.30 (https://
docs.stingray.science). The Z2 search and related calculations of confi-
dence level and chance probabilities are performed using Interactive
Data Language (IDL) version 8.7.3. The IDL codes developed can be
provided upon request. For spectral analysis we used the software
XSPEC, version v12.10.1f, (including the XANADU package) available
at https://heasarc.gsfc.nasa.gov/xanadu/xspec/.
30.	 Chanrion, O. et al. The Modular Multispectral Imaging Array (MMIA) of the ASIM payload
on the International Space Station. Space Sci. Rev. 215, 28 (2019).
31.
Vaughan, S. A Bayesian t
\end{document}
